\documentclass[11pt,a4paper]{article}

\usepackage[utf8]{inputenc} % allow utf-8 input
\usepackage[T1]{fontenc}    % use 8-bit T1 fonts
\usepackage{hyperref}       % hyperlinks
\usepackage{url}            % simple URL typesetting
\usepackage{booktabs}       % professional-quality tables
\usepackage{amsfonts}       % blackboard math symbols
\usepackage{nicefrac}       % compact symbols for 1/2, etc.
\usepackage{microtype}      % microtypography
\usepackage{xcolor}         % colors
\usepackage{graphicx}       % graphics
\usepackage[normalem]{ulem} % strikethrough
\title{Images in Markdown}

\begin{document}

\maketitle

A paragraph before the picture.

![A red box](data:image/png;base64,iVBORw0KGgoAAAANSUhEUgAAABgAAAAQCAIAAACDRijCAAAAHElEQVR4nGO8IyfHQA3ARBVTRg0aNWjUoBFsEABJ6QE4U1cKUgAAAABJRU5ErkJggg==)

An HTML \texttt{<img>} tag goes through the same policy:

\begin{figure}[h]
% image
\caption{A blue box}
\end{figure}

A relative path never resolves without the \texttt{local} tier: ![missing](img/not-here.png)

Trailing paragraph.

\end{document}